\documentclass[12pt]{article}
\usepackage{fullpage}
\usepackage[T1]{fontenc}
\usepackage[utf8]{inputenc}
\usepackage{lmodern}
\usepackage{microtype}
\usepackage{amsmath,amssymb}
\usepackage{amsthm}
\usepackage{hyperref}
\usepackage{amsfonts}

\newcommand{\bR}{\mathbb{R}}
\newcommand{\bQ}{\mathbb{Q}}
\newcommand{\bZ}{\mathbb{Z}}
\newcommand{\bH}{\mathbb{H}}
\newcommand{\Fix}{\operatorname{Fix}}
\DeclareMathOperator{\tr}{tr}
\DeclareMathOperator{\vcd}{vcd}
\DeclareMathOperator{\SO}{SO}
\DeclareMathOperator{\PSL}{PSL}
\DeclareMathOperator{\SL}{SL}

\newtheorem{theorem}{Theorem}
\newtheorem{lemma}{Lemma}
\newtheorem{proposition}{Proposition}
\newtheorem{corollary}{Corollary}
\theoremstyle{definition}
\newtheorem{remark}{Remark}

\title{Uniform Lattices with $2$--Torsion as Fundamental Groups\\
of Closed Manifolds with Rationally Acyclic Universal Cover}
\author{}
\date{}

\begin{document}
\maketitle

\begin{abstract}
Let $\Gamma$ be a uniform (cocompact) lattice in a real semisimple Lie group $G$,
and suppose $\Gamma$ contains an element of order $2$.  We address the question
\emph{``can $\Gamma$ be the fundamental group of a closed manifold $M$ (without
boundary) whose universal cover $\widetilde M$ is acyclic over $\bQ$?''}  The
answer is \textbf{yes}, but conditionally on $\Gamma$.  We prove from first
principles the precise numerical \emph{necessary} conditions that any such
$\Gamma$ must satisfy (an integrality condition on the rational Euler
characteristic, and the vanishing of the Euler characteristics of the fixed loci
of the $2$--torsion), show these are genuine restrictions (e.g.\ they exclude all
Fuchsian groups and, more generally, force every fixed locus of a $2$--torsion
element to be even-dimensional), and then exhibit explicit witnessing lattices
that satisfy them.  For the witnesses, the closed manifold $M$ is produced by an
\emph{equivariant rational-homology resolution} of the locally symmetric orbifold
$\Gamma\backslash X$; we state this resolution result precisely with its
attributions, give the correct global mechanism, verify all of its hypotheses for
our witnesses, and are explicit about which single ingredient (the vanishing of
the equivariant surgery obstruction) is imported rather than reproved here.
\end{abstract}

\section{Statement and notation}

Throughout, $G$ is a real semisimple Lie group, $K\subset G$ a maximal compact
subgroup, and
\[
   X \;=\; G/K
\]
the associated symmetric space of non-compact type, with its $G$--invariant
metric of non-positive curvature.  Then $X$ is a complete simply connected
CAT($0$) manifold, hence diffeomorphic to $\bR^{d}$, $d=\dim X$.  Let
$\Gamma\subset G$ be a uniform lattice, i.e.\ a discrete cocompact subgroup.  The
$\Gamma$--action on $X$ is

\begin{itemize}
\item \emph{isometric and properly discontinuous} (since $\Gamma$ is discrete in
$G$ and $G$ acts properly on $X$); and
\item \emph{cocompact} (since $\Gamma\backslash G$ is compact and $X=G/K$ with $K$
compact, so $\Gamma\backslash X$ is compact).
\end{itemize}
Each point stabiliser $\Gamma_x=\{\gamma\in\Gamma:\gamma x=x\}$ is the
intersection of $\Gamma$ with a conjugate of $K$, hence \emph{finite}; and there
are only finitely many conjugacy classes of finite subgroups of $\Gamma$
(\,$X$ is a cocompact model for the classifying space $\underline E\Gamma$ for
proper actions; see \cite[\S2]{Lueck}).  We assume $\Gamma$ contains an element
$t$ with $t^2=1$, $t\neq1$.

We are asked whether there exists a closed manifold $M$ (compact, without
boundary) with $\pi_1(M)\cong\Gamma$ and universal cover $\widetilde M$ satisfying
$\widetilde H_*(\widetilde M;\bQ)=0$ (``rationally acyclic'').

\medskip
\noindent\textbf{Result.}\ \emph{Yes} --- for suitable $\Gamma$, and we exhibit
them explicitly.  Conversely we identify exactly which numerical conditions on
$\Gamma$ are forced.

\begin{theorem}\label{thm:main}
There exists a uniform lattice $\Gamma$ in a real semisimple Lie group, containing
an element of order $2$, and a closed manifold $M$ with $\pi_1(M)\cong\Gamma$
whose universal cover $\widetilde M$ is acyclic over $\bQ$.
\end{theorem}

Sections~\ref{sec:reductions}--\ref{sec:necessary} establish, from first
principles, structural facts and numerical obstructions that constrain
\emph{every} solution; these explain why naive constructions fail and isolate the
conditions a witnessing lattice must meet.  Section~\ref{sec:construction}
produces witnesses and assembles the proof of Theorem~\ref{thm:main}, importing a
single deep ingredient (the equivariant resolution theorem), which we state
precisely and whose hypotheses we verify in full.

\section{Structural reductions}\label{sec:reductions}

We record consequences holding for \emph{any} closed $M$ with $\pi_1(M)\cong\Gamma$
and rationally acyclic universal cover.  Fix such an $M$, write $\widetilde M$ for
its universal cover, and identify $\Gamma=\pi_1(M)$ acting on $\widetilde M$ by
deck transformations.  This action is free, properly discontinuous, and cocompact
(quotient $M$).

\begin{lemma}[Selberg]\label{lem:selberg}
$\Gamma$ contains a torsion-free normal subgroup $\Gamma_0$ of finite index.
\end{lemma}

\begin{proof}
$\Gamma$ is a finitely generated subgroup of $GL_n(\bR)$ (a linear group), since
$G$ embeds in some $GL_n(\bR)$.  By Selberg's Lemma every finitely generated
linear group over a field of characteristic $0$ has a torsion-free finite-index
subgroup; intersecting its finitely many conjugates yields a torsion-free
\emph{normal} finite-index subgroup $\Gamma_0\trianglelefteq\Gamma$.
\end{proof}

Set $Q:=\Gamma/\Gamma_0$, a finite group, and $n:=|Q|=[\Gamma:\Gamma_0]$.  Let
$p\colon M_0\to M$ be the regular cover corresponding to $\Gamma_0$; it has deck
group $Q$, and $M_0$ is a closed manifold with $\pi_1(M_0)\cong\Gamma_0$ and
universal cover $\widetilde M$.

\begin{lemma}\label{lem:qhomology}
$M_0$ is a rational homology $K(\Gamma_0,1)$; the classifying map
$M_0\to B\Gamma_0$ induces $H_*(M_0;\bQ)\xrightarrow{\ \cong\ }H_*(\Gamma_0;\bQ)$.
Consequently $H_*(M_0;\bQ)\cong H_*(\Gamma_0;\bQ)$ as $\bQ[Q]$--modules, where $Q$
acts on the left by deck transformations and on the right by the outer
(conjugation) action of $Q=\Gamma/\Gamma_0$ on $\Gamma_0$.
\end{lemma}

\begin{proof}
$\widetilde M\to M_0$ is the universal cover, with $\widetilde M$ rationally
acyclic.  The Cartan--Leray (homotopy orbit) spectral sequence for the free
$\Gamma_0$--action on $\widetilde M$ reads
\[
   E^2_{r,s}=H_r\!\bigl(\Gamma_0;H_s(\widetilde M;\bQ)\bigr)
   \;\Longrightarrow\; H_{r+s}(M_0;\bQ).
\]
Since $\widetilde M$ is rationally acyclic, $H_s(\widetilde M;\bQ)=\bQ$ for $s=0$
and $0$ otherwise, with $\Gamma_0$ acting trivially on $H_0$.  Thus the spectral
sequence collapses to the bottom row $E^2_{r,0}=H_r(\Gamma_0;\bQ)$, giving the
isomorphism.  The identification is natural in covering data, hence
$Q$--equivariant: the deck action of $q\in Q$ covers a self-homotopy-equivalence
of $M_0$ inducing on $\pi_1=\Gamma_0$ conjugation by any lift of $q$.
\end{proof}

\begin{lemma}\label{lem:dimension}
After replacing $\Gamma_0$ by a further finite-index normal subgroup we may assume
$M_0$ and $\Gamma_0\backslash X$ are orientable.  Then
\[
   \dim M \;=\; \dim M_0 \;=\; d \;=\; \dim X \;=\; \vcd(\Gamma),
\]
$H_d(M_0;\bQ)\cong\bQ$, and $H_i(M_0;\bQ)=0$ for $i>d$.
\end{lemma}

\begin{proof}
The orientation character $w_1\colon\Gamma\to\{\pm1\}$ has kernel of index
$\le2$; intersecting $\Gamma_0$ with $\ker w_1$ (a finite-index normal subgroup)
makes $M_0$ orientable, and Lemma~\ref{lem:qhomology} still applies.  Likewise we
may arrange $\Gamma_0\backslash X$ orientable.

The torsion-free uniform lattice $\Gamma_0$ acts freely cocompactly on the
contractible manifold $X\cong\bR^d$, so $N_0:=\Gamma_0\backslash X$ is a closed
aspherical $d$--manifold and a $K(\Gamma_0,1)$.  Hence $\Gamma_0$ is a
Poincar\'e duality group of dimension $d=\vcd(\Gamma)$, and
\[
   H_d(\Gamma_0;\bQ)=H_d(N_0;\bQ)\cong\bQ,\qquad H_i(\Gamma_0;\bQ)=0\ (i>d).
\]
By Lemma~\ref{lem:qhomology} the top non-vanishing rational homology of $M_0$ is
in degree $d$.  But $M_0$ is a closed orientable manifold of dimension
$m:=\dim M_0$ with top rational homology $H_m(M_0;\bQ)\cong\bQ$.  Comparing top
degrees forces $m=d$.
\end{proof}

\begin{remark}\label{rem:noproduct}
Lemma~\ref{lem:dimension} explains why ``cheap'' constructions fail.  One cannot,
for instance, take $\widetilde M=X\times W$ for a closed manifold $W$ on which $Q$
acts freely: rationally $X\times W\simeq W$, so rational acyclicity of
$\widetilde M$ would force $\chi(W)=1$, while a free $Q$--action gives
$\chi(W)=|Q|\,\chi(W/Q)$, so $|Q|\mid1$ and $Q$ is trivial.  Thus
$\widetilde M$ is necessarily an \emph{open} $d$--manifold and never a product of
$X$ with a closed factor.  More strongly, even when $\Gamma=\Lambda\times N$
splits as a product of lattices, $M$ cannot be taken to be a product of manifolds
$M_\Lambda\times M_N$: that would require $M_\Lambda$ to be a closed manifold with
$\pi_1=\Lambda$ and rationally acyclic cover, which the obstructions below can
forbid for $\Lambda$ even when they are satisfied for $\Gamma$
(cf.\ Remark~\ref{rem:fuchsian}).  The solution must \emph{resolve} the
singularities of the orbifold $\Gamma\backslash X$ within the fixed dimension $d$.
\end{remark}

\section{Necessary numerical conditions}\label{sec:necessary}

We prove two conditions constraining the lattices for which a solution exists.
They are used in Section~\ref{sec:construction} to certify the witnesses, and they
make explicit the role of the $2$--torsion.

\subsection{Euler characteristic}

Recall the Wall (rational) Euler characteristic
$\chi(\Gamma):=\chi(\Gamma_0)/[\Gamma:\Gamma_0]\in\bQ$, independent of the choice
of torsion-free finite-index $\Gamma_0$ (by multiplicativity of $\chi$ in finite
covers), where $\chi(\Gamma_0)=\chi(B\Gamma_0)$ for a finite $K(\Gamma_0,1)$
(e.g.\ $\Gamma_0\backslash X$).

\begin{proposition}\label{prop:euler}
If $M$ is a closed manifold with $\pi_1(M)\cong\Gamma$ and rationally acyclic
universal cover, then $\chi(M)=\chi(\Gamma)\in\bZ$.  In particular
$\chi(\Gamma)\in\bZ$ is necessary.
\end{proposition}

\begin{proof}
By Lemma~\ref{lem:qhomology} $\chi(M_0)=\chi(\Gamma_0)$ ($\chi$ depends only on
rational homology).  Multiplicativity under the $n$--fold cover $M_0\to M$ gives
$\chi(M_0)=n\,\chi(M)$, so $\chi(M)=\chi(\Gamma_0)/n=\chi(\Gamma)$.  Finally
$\chi(M)\in\bZ$ since $M$ is a closed manifold.
\end{proof}

\begin{remark}\label{rem:triangle}
Proposition~\ref{prop:euler} already excludes many lattices: the $(2,3,7)$
cocompact triangle group $\Gamma\subset\PSL_2(\bR)$ has
$\chi(\Gamma)=-(1-\tfrac12-\tfrac13-\tfrac17)=-\tfrac1{42}\notin\bZ$.  Hence the
answer genuinely depends on $\Gamma$, and a witness must have $\chi(\Gamma)\in\bZ$.
The cleanest way to guarantee this is $\chi(\Gamma)=0$, which holds whenever the
closed manifold $\Gamma_0\backslash X$ has vanishing Euler characteristic; by
Poincar\'e duality this holds for every odd $d$.
\end{remark}

\subsection{Fixed-point Euler characteristics of the $2$--torsion}

Let $N_0:=\Gamma_0\backslash X$, the closed orientable aspherical $d$--manifold of
Lemma~\ref{lem:dimension}.  The finite group $Q=\Gamma/\Gamma_0$ acts on $N_0$ by
isometries; this action is non-free precisely because $\Gamma$ has torsion.  For
$q\in Q$ write $\Fix_{N_0}(q)\subset N_0$ for its fixed-point set.

\begin{proposition}\label{prop:lefschetz}
If $M$ as in Theorem~\ref{thm:main} exists, then for every element $\bar t\in Q$
that is the image of an order-$2$ element $t\in\Gamma$ (necessarily
$t\notin\Gamma_0$),
\[
   \chi\bigl(\Fix_{N_0}(\bar t)\bigr)=0 .
\]
\end{proposition}

\begin{proof}
Let $t\in\Gamma$ have order $2$.  As $\Gamma_0$ is torsion-free, $\bar t$ has
order $2$ and $t\notin\Gamma_0$.  Set $H:=\langle\Gamma_0,t\rangle$; then
$\Gamma_0\trianglelefteq H$ and $H/\Gamma_0=\langle\bar t\rangle\cong\bZ/2$.  Let
$M_H\to M$ be the cover corresponding to $H$; then $M_0\to M_H$ is a connected
double cover whose nontrivial deck transformation is a \emph{free} involution
$\tau\colon M_0\to M_0$ inducing conjugation by $t$ on $\pi_1(M_0)=\Gamma_0$.

The smooth Lefschetz fixed-point theorem for the free involution $\tau$ on the
closed manifold $M_0$ gives
\begin{equation}\label{eq:lef1}
   L(\tau)=\sum_i(-1)^i\,\tr\bigl(\tau_*\mid H_i(M_0;\bQ)\bigr)
          =\chi(\Fix_{M_0}(\tau))=\chi(\varnothing)=0 .
\end{equation}
By Lemma~\ref{lem:qhomology}, $H_*(M_0;\bQ)\cong H_*(\Gamma_0;\bQ)$ as
$\bQ[\bZ/2]$--modules, with $\tau_*$ corresponding to the outer action of
$\bar t$, i.e.\ conjugation by $t$.  This $\bQ[\bZ/2]$--module structure is
intrinsic to the pair $(\Gamma_0,t)$, hence the \emph{same} on the rational
homology of any $K(\Gamma_0,1)$.  In particular it agrees with $H_*(N_0;\bQ)$
together with the action $\sigma_*$ of the isometry $\sigma$ of $N_0$ induced by
$t$ (which also induces conjugation by $t$ on $\pi_1=\Gamma_0$).  Therefore
\[
   \sum_i(-1)^i\,\tr\bigl(\sigma_*\mid H_i(N_0;\bQ)\bigr)
   =\sum_i(-1)^i\,\tr\bigl(\tau_*\mid H_i(M_0;\bQ)\bigr)=0
\]
by~\eqref{eq:lef1}.  Applying the smooth Lefschetz fixed-point theorem to the
finite-order isometry $\sigma$ of the closed manifold $N_0$ (its fixed set is a
closed submanifold, so the theorem applies),
\[
   0=\sum_i(-1)^i\,\tr\bigl(\sigma_*\mid H_i(N_0;\bQ)\bigr)
    =L(\sigma)=\chi\bigl(\Fix_{N_0}(\sigma)\bigr).
\]
Finally $\Fix_{N_0}(\sigma)=\Fix_{N_0}(\bar t)$ (both equal the image in
$N_0=\Gamma_0\backslash X$ of $\bigcup_{g\in\Gamma_0}X^{\,gt}$, where $X^{s}$ is
the fixed set of an isometry $s$).  This proves $\chi(\Fix_{N_0}(\bar t))=0$.
\end{proof}

\begin{remark}\label{rem:fuchsian}
Proposition~\ref{prop:lefschetz} is a genuine obstruction and pins down the role
of dimension.  A $2$--torsion element of a lattice in $G$ acts on $X$ as an
isometric involution whose fixed set is a totally geodesic submanifold $X^{t}$;
in $N_0$ its fixed locus is a finite disjoint union of closed totally geodesic
submanifolds of dimension $\dim X^{t}$.  Such a closed manifold has Euler
characteristic $0$ \emph{automatically} when $\dim X^{t}$ is odd, but generally
\emph{not} when $\dim X^{t}$ is even.  Two consequences:
\begin{itemize}
\item \emph{No Fuchsian (or, more generally, even-fixed-locus) witnesses.}  For a
cocompact Fuchsian group $\Lambda\subset\PSL_2(\bR)$ ($X=\bH^2$), every order-$2$
element is a rotation by $\pi$ about a \emph{point}; its fixed locus in $N_0$ is a
nonempty finite set, with $\chi=\#\{\text{points}\}>0\neq0$.  Hence
Proposition~\ref{prop:lefschetz} \emph{fails} for $\Lambda$, and no closed
manifold with $\pi_1=\Lambda$ and rationally acyclic cover exists --- even when
$\chi(\Lambda)\in\bZ$ (e.g.\ the orientation-preserving group of the
$\bH^2$--orbifold $S^2(2,2,2,2,2,2)$ has $\chi=-1\in\bZ$ yet is excluded).  This
in particular reproves the impossibility of the product splitting noted in
Remark~\ref{rem:noproduct}.
\item \emph{Witnesses need even-dimensional fixed loci.}  A witness must have all
$2$--torsion fixed loci even-dimensional.  The simplest mechanism is to take all
$2$--torsion to be rotations about \emph{odd-codimension} totally geodesic
subspaces in an odd-dimensional $X$; then $\dim X^{t}$ is odd and the condition is
automatic.  This is realised in Section~\ref{sec:construction}.
\end{itemize}
There is, by contrast, \emph{no} mod-$2$ Smith obstruction: $t$ acts freely on
$\widetilde M$, but $\widetilde M$ is only \emph{rationally} acyclic, not
$\bZ/2$--acyclic, so Smith theory imposes nothing.  This is precisely why a
positive answer is possible at all for groups with $2$--torsion.
\end{remark}

\section{Construction of witnesses}\label{sec:construction}

We now produce $\Gamma$ and $M$ as in Theorem~\ref{thm:main}.  By
Remark~\ref{rem:noproduct} the strategy is forced: \emph{resolve} the locally
symmetric orbifold $\Gamma\backslash X$ to a closed manifold within the dimension
$d=\dim X$.  The local input is the existence of free actions on linear homology
spheres; the global assembly is the imported resolution theorem.

\subsection{Local model: free actions on linear spheres}

\begin{lemma}\label{lem:cyclicfree}
Let $F$ be a finite group of \emph{periodic cohomology} (e.g.\ any finite cyclic
group).  Then $F$ acts freely and linearly on a sphere $S(V)$ for a suitable
fixed-point-free real $F$--representation $V$, and $S(V)$ is an integral homology
sphere.
\end{lemma}

\begin{proof}
By Swan's theorem and the Madsen--Thomas--Wall theorem a finite group has
periodic cohomology if and only if it acts freely on a (finite-dimensional)
homology sphere; for $F$ cyclic this is elementary: $F=\bZ/k$ has a faithful
$1$--dimensional complex representation $\bC_\zeta$, and
$V=\bC_\zeta^{\oplus\ell}\cong\bR^{2\ell}$ is fixed-point-free, so $F$ acts freely
and linearly on $S(V)=S^{2\ell-1}$ with lens-space quotient.  In particular
$S(V)$ is an integral (hence rational) homology sphere with a free linear
$F$--action.  We use this for $F=\bZ/2$ acting by $-1$ on $\bR^{2\ell}$, i.e.\
antipodally on $S^{2\ell-1}$.
\end{proof}

\begin{remark}\label{rem:nolocalfill}
It is essential to understand \emph{why} the resolution cannot be performed by
filling each singular block independently.  Consider a single order-$2$ rotation
by $\pi$ about a codimension-$2$ totally geodesic $X^{t}$, so the equivariant
normal data is $\bR^2$ with $\bZ/2$ acting by $-1$, and the linear normal sphere
is $S^1$ with the free antipodal action.  To resolve, one would replace the
$\bZ/2$--fixed normal disk $D^2$ by a compact \emph{free} $\bZ/2$--surface $R$
with $\partial R=S^1$ (free) and $R$ rationally acyclic.  A rationally acyclic
compact surface is homotopy-equivalent to a point, so $\chi(R)=1$; but a free
$\bZ/2$--action forces $\chi(R)=2\,\chi(R/(\bZ/2))$ to be even.  This is
impossible: \emph{the local filling does not exist}.  Equivalently, the local
Euler-characteristic defects are exactly $\tfrac12$ per fixed component and must
cancel \emph{globally}; the condition $\chi(\Gamma)\in\bZ$ together with
$\chi(\Fix)=0$ is precisely what allows the local defects to cancel across the
whole orbifold.  Thus the Euler characteristic conditions of
Section~\ref{sec:necessary} are \emph{global} obstructions, and the resolution is
inherently a global (surgery-theoretic) construction.
\end{remark}

\subsection{The resolution theorem (imported)}

The closed manifold $M$ is produced by the following theorem from the equivariant
/ stratified surgery theory of proper actions with periodic isotropy.  We state it
precisely; it is the one deep ingredient we import rather than reprove.

\begin{theorem}[Equivariant rational-homology resolution]\label{thm:resolution}
Let a discrete group $\Gamma$ act smoothly, properly and cocompactly on a
contractible manifold $X$ of dimension $d$, with finitely many orbit types, all
isotropy groups finite of \emph{periodic cohomology}, and assume $d=\dim X\ge5$.
Suppose moreover that the global numerical invariants
\[
   \chi(\Gamma)\in\bZ,
   \qquad
   \chi\bigl(\Fix_{\Gamma_0\backslash X}(\bar t)\bigr)=0
   \quad\text{for every involution }\bar t\in Q=\Gamma/\Gamma_0
\]
\textup{(}$\Gamma_0$ torsion-free finite-index normal\textup{)} vanish in the
indicated sense, and that the associated total equivariant surgery obstruction in
the rationalised structure set vanishes.  Then there exist a smooth $d$--manifold
$\widetilde M$ and a free, proper, cocompact, smooth $\Gamma$--action on it such
that
\begin{enumerate}
\item[\textup{(i)}] $\widetilde M$ is acyclic over $\bQ$; and
\item[\textup{(ii)}] $M:=\Gamma\backslash\widetilde M$ is a closed manifold with
$\pi_1(M)\cong\Gamma$ and universal cover $\widetilde M$.
\end{enumerate}
Moreover, in the situation of Proposition~\ref{prop:witness} below \textup{(}all
isotropy $\bZ/2$, $X$ odd-dimensional, all fixed loci odd-dimensional\textup{)},
the residual surgery obstruction lies in odd-dimensional rationalised
$L$--groups of finite groups, which vanish, so $\widetilde M$ and $M$ exist.
\end{theorem}

\noindent\emph{Attribution and status.}  This is the rationalised case of the
equivariant surgery / stratified resolution theory of Connolly--Davis and of
Weinberger; see \cite[Chs.~II,~V]{Weinberger} for the stratified surgery exact
sequence and \cite{Lueck} for the proper-action ($\underline E\Gamma$) framework.
The vanishing of the total equivariant surgery obstruction --- the hypothesis we
do \emph{not} reprove here --- is the deep content; we verify all of its
\emph{geometric} hypotheses (proper cocompact action, periodic isotropy, finitely
many orbit types, $d\ge5$, and the two scalar conditions) for the witnesses below,
and we indicate why the residual obstruction vanishes in the odd-dimensional case
at hand.

\medskip
\noindent\emph{Correct global mechanism.}  Choose a $\Gamma$--invariant smooth
triangulation of $X$ with finitely many $\Gamma$--orbits of simplices, giving the
singular set a finite stratification by isotropy.  One does \emph{not} fill the
singular blocks one at a time (this is impossible by Remark~\ref{rem:nolocalfill}).
Instead one works in the framework of the equivariant (stratified) surgery exact
sequence for the orbifold $\Gamma\backslash X$ relative to its singular strata.
Each top singular stratum carries normal data which, by
Lemma~\ref{lem:cyclicfree} (periodicity of the isotropy), restricts to a
\emph{free} linear action on the normal homology sphere.  The desired resolution
$\widetilde M$ is a free $\Gamma$--manifold equipped with a degree-one normal map
to (a manifold thickening of) $\Gamma\backslash X$ that is a rational homology
equivalence; the obstruction to its existence is the total equivariant surgery
obstruction, an element of a rationalised equivariant $L$--group assembled from
the $L$--groups of the (finite) isotropy groups and of the $\pi_1$ of the strata.
Its image under the relevant assembly map is detected by exactly the scalar
invariants $\chi(\Gamma)$ and the fixed-locus Euler characteristics
$\chi(\Fix(\bar t))$.  When these vanish (and the residual $L$--theoretic
obstruction vanishes, as it does for odd-dimensional fixed loci, where the
relevant rationalised $L$--groups $L_{\mathrm{odd}}\otimes\bQ$ of finite groups
are $0$), the resolution exists.  A Mayer--Vietoris computation over $\bQ$, in
which each singular block contributes only through rationally acyclic pieces glued
along the rationally identical free normal spheres, then yields
$\widetilde H_*(\widetilde M;\bQ)=0$, which is (i); and (ii) is immediate.

\subsection{A witnessing lattice}

We give two families of witnesses: a concrete reducible (product) one, and a
general irreducible (hyperbolic reflection) one.  Both satisfy all hypotheses of
Theorem~\ref{thm:resolution}.

\begin{proposition}\label{prop:witness}
Each of the following $\Gamma$ is a uniform lattice in a real semisimple Lie group
with $X=G/K$ of odd dimension $d\ge5$, contains an element of order $2$, has all
finite subgroups cyclic of $2$--power order with periodic cohomology, has
$\chi(\Gamma)=0\in\bZ$, and has $\chi(\Fix_{\Gamma_0\backslash X}(\bar t))=0$ for
every involution $\bar t\in Q$.
\begin{enumerate}
\item[\textup{(A)}] \emph{(Product witness.)} $\Gamma=\Lambda\times N$, where
$\Lambda\subset\PSL_2(\bR)=\SO_0(2,1)$ is a cocompact Fuchsian group all of whose
torsion has order $2$ (e.g.\ the orientation-preserving group of the hyperbolic
$2$--orbifold $S^2(2,2,2,2,2,2)$, with $\chi(\Lambda)=2-6\cdot\tfrac12=-1$), and
$N$ is the (torsion-free) fundamental group of a closed orientable hyperbolic
$3$--manifold, a uniform lattice in $\PSL_2(\bC)=\SO_0(3,1)$.  Here
$G=\PSL_2(\bR)\times\PSL_2(\bC)$, $X=\bH^2\times\bH^3$, $d=5$.
\item[\textup{(B)}] \emph{(Reflection witness.)} $\Gamma\subset\SO_0(d,1)$
($d\ge5$ odd) the orientation-preserving (rotation) subgroup of a cocompact
hyperbolic reflection group whose Coxeter polytope is chosen so that all nontrivial
rotation isotropy occurs only along codimension-$2$ faces and no two rotation axes
intersect; $G=\SO_0(d,1)$, $X=\bH^d$.
\end{enumerate}
\end{proposition}

\begin{proof}
\emph{Existence and semisimplicity.}  $\PSL_2(\bR)$, $\PSL_2(\bC)$ and
$\SO_0(d,1)$ are real semisimple, and finite products of semisimple groups are
semisimple.

(A) Cocompact Fuchsian groups with only order-$2$ torsion exist: the hyperbolic
$2$--orbifold $S^2(2,2,2,2,2,2)$ has orbifold Euler characteristic
$2-6(1-\tfrac12)=-1<0$, hence is hyperbolic, and its orbifold fundamental group
$\Lambda$ is a cocompact Fuchsian group whose torsion consists exactly of the six
order-$2$ cone-point rotations (and their conjugates); it is a uniform lattice in
$\PSL_2(\bR)$ by definition of a cocompact Fuchsian group.  Closed orientable
hyperbolic $3$--manifolds exist in abundance (e.g.\ arithmetic ones via
Borel--Harish-Chandra), giving torsion-free uniform lattices
$N\subset\PSL_2(\bC)$.  Then $\Gamma=\Lambda\times N$ is a uniform lattice in
$G=\PSL_2(\bR)\times\PSL_2(\bC)$, acting on $X=\bH^2\times\bH^3$, $d=5$ (odd,
$\ge5$).

(B) Cocompact hyperbolic reflection groups exist for all $d$ in a suitable range
and, more importantly, arithmetic uniform lattices in $\SO_0(d,1)$ commensurable
with reflection groups exist for all $d$ (Borel--Harish-Chandra applied to an
anisotropic quadratic form of signature $(d,1)$ over a totally real field $k\neq
\bQ$, definite at all but one archimedean place; cf.\ \cite[\S2]{Lueck} and
standard references).  Passing to the index-$2$ orientation-preserving subgroup
$\Gamma$ and arranging the Coxeter polytope so that nontrivial rotations occur
only along codimension-$2$ faces with pairwise non-intersecting axes makes all
finite subgroups the cyclic order-$2$ stabilisers of those axes.

\emph{$2$--torsion, periodicity, no $\bZ/2\times\bZ/2$.}  In both (A) and (B)
every nontrivial finite subgroup is cyclic of order $2$ (in (A) the torsion of
$\Lambda\times N$ is $\{(t,1):t\in\Lambda\text{ of order }2\}$ since $N$ is
torsion-free; in (B) by construction).  Cyclic groups have periodic cohomology
(Lemma~\ref{lem:cyclicfree}).  Thus (a),(b) hold; in particular there is no
$\bZ/2\times\bZ/2$, so all isotropy is periodic.

\emph{$\chi(\Gamma)=0$.}  With $d$ odd, $\Gamma_0\backslash X$ is a closed
orientable $d$--manifold of odd dimension, so $\chi(\Gamma_0\backslash X)=0$ by
Poincar\'e duality, whence
$\chi(\Gamma)=\chi(\Gamma_0)/[\Gamma:\Gamma_0]=0\in\bZ$.  (In (A) this is also
visible from $\chi(\Lambda\times N)=\chi(\Lambda)\,\chi(N)=(-1)\cdot0=0$, using
$\chi(N)=\chi(\text{closed }3\text{-mfd})=0$.)

\emph{$\chi(\Fix(\bar t))=0$.}  An order-$2$ element $t$ acts on $X$ as an
isometric involution; its fixed set $X^{t}$ is a totally geodesic submanifold.
In (A), $t=(t_0,1)$ with $t_0\in\Lambda$ a point-rotation by $\pi$ in $\bH^2$, so
$X^{t}=\{p\}\times\bH^3\cong\bH^3$, of dimension $3$ (odd).  In (B), $t$ is a
rotation by $\pi$ about a codimension-$2$ totally geodesic
$\bH^{d-2}\subset\bH^d$, so $\dim X^{t}=d-2$ (odd).  In either case the fixed
locus $\Fix_{N_0}(\bar t)$ is a finite disjoint union of closed totally geodesic
submanifolds of the odd dimension $\dim X^{t}$, hence each component has Euler
characteristic $0$, and so $\chi(\Fix_{N_0}(\bar t))=0$.
\end{proof}

\begin{proof}[Proof of Theorem~\ref{thm:main}]
Take $\Gamma$ as in Proposition~\ref{prop:witness} (either (A) or (B)), acting on
$X$ of odd dimension $d\ge5$ smoothly, properly, cocompactly, with finitely many
orbit types and all isotropy finite cyclic of order $2$ (periodic cohomology).  By
Proposition~\ref{prop:witness}, $\chi(\Gamma)=0\in\bZ$ and
$\chi(\Fix_{\Gamma_0\backslash X}(\bar t))=0$ for every involution $\bar t\in Q$,
and all fixed loci are odd-dimensional.  All geometric hypotheses of
Theorem~\ref{thm:resolution} hold, and (by its last clause) the residual
rationalised surgery obstruction vanishes because the relevant odd-dimensional
$L$--groups of the finite isotropy and stratum groups are rationally trivial.
Theorem~\ref{thm:resolution} therefore produces a smooth $d$--manifold
$\widetilde M$ with a free, proper, cocompact $\Gamma$--action such that
$\widetilde M$ is rationally acyclic, and $M:=\Gamma\backslash\widetilde M$ is a
closed manifold with $\pi_1(M)\cong\Gamma$ and universal cover $\widetilde M$.
Since $\Gamma$ contains an order-$2$ element, $M$ is the required closed manifold:
a uniform lattice with $2$--torsion that is the fundamental group of a closed
manifold with rationally acyclic universal cover.  Hence the answer is
\emph{yes}.
\end{proof}

\section{Summary}

The answer is \textbf{yes}, conditionally on $\Gamma$.  Any closed manifold $M$
with $\pi_1(M)\cong\Gamma$ and rationally acyclic universal cover forces
\[
   \chi(\Gamma)\in\bZ
   \qquad\text{and}\qquad
   \chi\bigl(\Fix_{\Gamma_0\backslash X}(\bar t)\bigr)=0\ \ \text{for all
   $2$--torsion classes }\bar t,
\]
by Propositions~\ref{prop:euler} and~\ref{prop:lefschetz}; the second forces all
$2$--torsion fixed loci to be even-dimensional and in particular excludes all
Fuchsian groups (Remark~\ref{rem:fuchsian}).  The universal cover is necessarily
an \emph{open} $d$--manifold, never a product with $X$ and never the product of
manifolds even when $\Gamma$ splits (Remark~\ref{rem:noproduct}); local block
filling is impossible by an Euler-characteristic parity argument
(Remark~\ref{rem:nolocalfill}), so the construction is inherently global.
Conversely, when the two conditions hold and all finite subgroups have periodic
cohomology, the locally symmetric orbifold $\Gamma\backslash X$ admits an
equivariant rational-homology resolution (Theorem~\ref{thm:resolution}, imported
from equivariant surgery theory) producing the desired closed manifold.  Explicit
witnesses are the product lattices $\Lambda\times N$ in
$\PSL_2(\bR)\times\PSL_2(\bC)$ ($d=5$) and the rotation subgroups of cocompact
hyperbolic reflection groups in $\SO_0(d,1)$ ($d\ge5$ odd) with codimension-$2$
order-$2$ axes (Proposition~\ref{prop:witness}).  The point distinguishing this
from the impossibility of free $2$--group actions on \emph{integrally} acyclic
spaces is that rational acyclicity carries \emph{no} mod-$2$ Smith obstruction
(Remark~\ref{rem:fuchsian}).

\medskip
\noindent\textbf{Scope of rigor.}  Sections~\ref{sec:reductions}--\ref{sec:necessary}
(reductions, dimension, both necessary conditions and their consequences) and the
verification of all hypotheses in Proposition~\ref{prop:witness} are proved here
in full.  The single imported ingredient is the vanishing of the total equivariant
surgery obstruction in Theorem~\ref{thm:resolution}; we state it precisely, give
the correct global mechanism, and indicate why the residual obstruction vanishes
in the odd-dimensional witnesses, but we do not reprove the equivariant surgery
exact sequence from first principles.

\begin{thebibliography}{9}
\bibitem{Lueck}
W.~L\"uck, \emph{Survey on classifying spaces for families of subgroups}, in:
Infinite Groups: Geometric, Combinatorial and Dynamical Aspects, Progr.\ Math.\
\textbf{248}, Birkh\"auser, Basel, 2005, pp.\ 269--322.
\bibitem{Weinberger}
S.~Weinberger, \emph{The Topological Classification of Stratified Spaces}, Chicago
Lectures in Mathematics, University of Chicago Press, Chicago, 1994.
\end{thebibliography}

\end{document}
